% Bien dich bang XeLaTeX hai lan de cap nhat muc luc va tham chieu.
% File doc lap: khong can anh, BibTeX hay cac file Python khi bien dich.
\documentclass[12pt,a4paper,oneside]{report}
\usepackage{fontspec}
\usepackage{polyglossia}
\setdefaultlanguage{vietnamese}
\setmainfont{TeX Gyre Termes}
\setsansfont{TeX Gyre Heros}
\setmonofont{Latin Modern Mono}[Scale=MatchLowercase]
\usepackage[left=3cm,right=2.5cm,top=2.5cm,bottom=2.5cm]{geometry}
\usepackage{amsmath,amssymb,booktabs,longtable,tabularx,array}
\usepackage{graphicx,xcolor,listings,enumitem,setspace,fancyhdr}
\usepackage{tikz}
\usetikzlibrary{arrows.meta,positioning}
\usepackage{hyperref}
\hypersetup{colorlinks=true,linkcolor=blue!45!black,urlcolor=blue!60!black,
  citecolor=blue!45!black,pdftitle={Bao cao trien khai mo hinh GNN-DRL cho CVRP}}
\onehalfspacing
\setlength{\headheight}{16pt}
\pagestyle{fancy}
\fancyhf{}
\fancyhead[L]{\small Mô hình GNN--DRL cho CVRP}
\fancyhead[R]{\small Báo cáo triển khai}
\fancyfoot[C]{\thepage}
\setlist{itemsep=3pt,topsep=5pt}
\setlength{\emergencystretch}{3em}
\definecolor{codebg}{RGB}{247,248,250}
\definecolor{codeblue}{RGB}{30,70,150}
\lstdefinestyle{pythoncode}{language=Python,basicstyle=\ttfamily\footnotesize,
  backgroundcolor=\color{codebg},keywordstyle=\color{codeblue}\bfseries,
  commentstyle=\color{gray},stringstyle=\color{teal!70!black},
  numbers=left,numberstyle=\tiny\color{gray},numbersep=7pt,
  frame=single,breaklines=true,breakatwhitespace=false,
  columns=fullflexible,keepspaces=true,showstringspaces=false,tabsize=4,
  captionpos=b}
\lstset{style=pythoncode}
\newcommand{\file}[1]{\texttt{\detokenize{#1}}}
\newcommand{\E}{\mathbb{E}}
\newcommand{\ind}{\mathbf{1}}
\newcolumntype{Y}{>{\raggedright\arraybackslash}X}

\begin{document}
\begin{titlepage}
\centering
\vspace*{1.5cm}
{\Large\bfseries BÁO CÁO TRIỂN KHAI MÔ HÌNH\par}
\vspace{1.2cm}
{\LARGE\bfseries HỌC TĂNG CƯỜNG SÂU KẾT HỢP\par}
\vspace{0.3cm}
{\LARGE\bfseries MẠNG NƠ-RON ĐỒ THỊ CHO\par}
\vspace{0.3cm}
{\LARGE\bfseries BÀI TOÁN ĐỊNH TUYẾN XE\par}
\vspace{0.6cm}
{\large Với bất định nhu cầu và giao thông\par}
\vspace{1.4cm}
\begin{minipage}{0.9\textwidth}
\textbf{Đối tượng báo cáo:} Chương trình Python tại thư mục
\file{research/cvrp_reproduction}.\\[0.4cm]
\textbf{Cơ sở nghiên cứu:} Bài báo của Kadyrov và cộng sự (2025),
\emph{Deep reinforcement learning for dynamic vehicle routing with demand
and traffic uncertainty}.\\[0.4cm]
\textbf{Phạm vi:} Mô hình toán, thiết kế môi trường, kiến trúc mạng,
huấn luyện PPO, phương pháp đối chứng, đánh giá và kiểm thử mã nguồn.
\end{minipage}
\vfill
{\large Tháng 11 năm 2026\par}
\end{titlepage}

\pagenumbering{roman}
\chapter*{Tóm tắt}
\addcontentsline{toc}{chapter}{Tóm tắt}
Báo cáo trình bày chương trình tái dựng mô hình định tuyến một xe có giới
hạn sức chứa trong môi trường có nhu cầu khách hàng chưa biết và chi phí
di chuyển biến động. Chương trình sử dụng hai lớp NNConv để xử lý đồ thị,
Transformer để tăng cường tương tác giữa các nút, GRU để duy trì bộ nhớ và
Proximal Policy Optimization (PPO) để huấn luyện chính sách chọn điểm đến.
Ba phương pháp suy luận gồm GNN--DRL, GNN--DRL--G và GNN--DRL--CW được
đánh giá cùng bảy phương pháp đối chứng trên một môi trường thống nhất.

Nội dung tập trung giải thích các lớp, hàm và luồng dữ liệu đang được cài
đặt; đồng thời công khai những lựa chọn bổ sung ở các điểm bài báo chưa
mô tả đủ chi tiết. Theo biên bản kiểm tra triển khai, 69 kiểm thử tự động
đã đạt và các bước huấn luyện thử, lưu, nạp, tiếp tục checkpoint đã chạy
thành công. Các kết quả này xác nhận khả năng vận hành của phần mềm;
chưa xác nhận việc tái lập các bảng kết quả nghiên cứu sau huấn luyện dài.

\noindent\textbf{Từ khóa:} CVRP, GNN, NNConv, Transformer, GRU, PPO,
nhu cầu ngẫu nhiên, giao thông biến động.
\clearpage
\tableofcontents
\clearpage
\listoffigures
\listoftables
\clearpage
\pagenumbering{arabic}

\chapter{Mục tiêu và tổ chức chương trình}
\section{Mục tiêu triển khai}
Chương trình chuyển ý tưởng của bài báo gốc~\cite{kadyrov2025} thành một
hệ thống nghiên cứu có thể huấn luyện, suy luận và đánh giá từ dòng lệnh.
Hệ thống giải quyết bài toán một xe xuất phát từ kho, phục vụ tập khách
hàng và quay về kho, cho phép nhiều chuyến lấy hàng. Các vị trí khách
hàng được xác định từ đầu; chương trình hiện chưa mô phỏng khách hàng
mới xuất hiện trong khi xe đang chạy.

Ba yêu cầu triển khai chính là: bảo toàn thông tin chưa quan sát được,
hạch toán đầy đủ chi phí mỗi lần di chuyển và báo cáo rõ các lượt chưa
phục vụ hết khách. Mô hình chỉ nhận nhu cầu danh nghĩa trước lần thăm đầu
tiên; biến chứa nhu cầu thật không được đưa trực tiếp vào quan sát.

\section{Cấu trúc mã nguồn}
\begin{table}[htbp]
\centering\small
\caption{Các thành phần chính của chương trình}
\begin{tabularx}{\textwidth}{@{}p{4.1cm}Y@{}}
\toprule
Tệp & Chức năng \\
\midrule
\file{cvrp/environment.py} & Cấu hình, sinh bài toán, quan sát, mask và chuyển trạng thái.\\
\file{cvrp/model.py} & Mạng RoutingAgent, tính NNConv theo từng nhóm cạnh, actor và critic.\\
\file{cvrp/train.py} & Thu thập quỹ đạo, PPO, log huấn luyện và checkpoint.\\
\file{cvrp/solvers.py} & Bảy baseline, hai mask lai và hàm giải chung.\\
\file{cvrp/evaluate.py} & Thực nghiệm ghép cặp, tổng hợp CSV và Friedman test.\\
\file{tests/} & Kiểm tra môi trường, mạng, gradient, huấn luyện và thống kê.\\
\file{reference/} & Notebook tham chiếu của tác giả, không tự động thực thi.\\
\bottomrule
\end{tabularx}
\end{table}

\section{Luồng hoạt động tổng thể}
Môi trường tạo đồ thị và quan sát hiện tại. Mạng sinh logits cho từng nút
và giá trị trạng thái. Sau khi loại các hành động không hợp lệ, chương
trình lấy mẫu điểm đến, thực thi trong môi trường và lưu kết quả. Trong
huấn luyện, quỹ đạo được dùng để cập nhật PPO. Trong đánh giá, trọng số
được giữ cố định và chỉ ghi chi phí, thời gian, mức hoàn thành.

\begin{figure}[htbp]
\centering
\begin{tikzpicture}[node distance=0.55cm,
  box/.style={draw=blue!55!black,fill=blue!4,rounded corners,
    text width=10.5cm,align=center,minimum height=0.85cm,font=\small},
  arrow/.style={-{Latex},thick}]
\node[box] (env) {CVRPEnv: vị trí, nhu cầu đã biết, tải và chi phí cạnh};
\node[box,below=of env] (gnn) {Hai lớp NNConv và Transformer};
\node[box,below=of gnn] (gru) {Global mean pooling và GRU};
\node[box,below=of gru] (heads) {Actor: logits từng nút \qquad Critic: giá trị trạng thái};
\node[box,below=of heads] (mask) {Mask khả thi hoặc mask lai heuristic};
\node[box,below=of mask] (step) {Lấy mẫu hành động, cập nhật môi trường, ghi reward};
\node[box,below=of step] (ppo) {Huấn luyện: cập nhật PPO \qquad Đánh giá: ghi chỉ số};
\draw[arrow] (env)--(gnn);
\draw[arrow] (gnn)--(gru);
\draw[arrow] (gru)--(heads);
\draw[arrow] (heads)--(mask);
\draw[arrow] (mask)--(step);
\draw[arrow] (step)--(ppo);
\end{tikzpicture}
\caption{Luồng dữ liệu của bản triển khai}
\end{figure}

\chapter{Mô hình bài toán và môi trường mô phỏng}
\section{Ký hiệu và mục tiêu}
Gọi $n$ là cạnh lưới, $M=n^2$ là số nút biểu diễn và
$\mathcal V=\{0,\ldots,M-1\}$. Nút $0$ là kho, tập
$\mathcal U\subset\mathcal V\setminus\{0\}$ chứa $N$ khách hàng.
Các ô không có khách vẫn được đưa vào mạng nhưng không được chọn làm
điểm giao hàng. Đây là cách biểu diễn theo notebook tham chiếu~\cite{notebook},
khác cách mô tả đồ thị chỉ có $N+1$ nút trong bài báo.

Tại bước $t$, xe ở $p_t$, có tải $L_t$ và chọn hành động $a_t$.
Mục tiêu định tuyến là giảm chi phí kỳ vọng:
\begin{equation}
\min_{\pi}\E_{\pi}\left[\sum_{t=0}^{T-1}c_{p_t,a_t}^{(t)}\right].
\end{equation}
Xe phải phục vụ đủ khách và kết thúc tại kho. Điều kiện sức chứa được
kiểm tra với nhu cầu thực khi đến khách; chương trình có quy tắc xử lý
thiếu hàng được trình bày tại mục~\ref{sec:recourse}.

\section{Lớp cấu hình EnvConfig}
\begin{table}[htbp]
\centering
\caption{Thông số môi trường mặc định}
\begin{tabularx}{\textwidth}{@{}lclY@{}}
\toprule
Tên trong mã & Ký hiệu & Mặc định & Ý nghĩa \\
\midrule
\file{grid_size} & $n$ & 8 & Lưới $8\times8$.\\
\file{customers} & $N$ & 20 & Số khách khi khởi tạo cấu hình thông thường.\\
\file{capacity} & $C$ & 15 & Sức chứa xe.\\
\file{beta} & $\beta$ & 0,07 & Hệ số giữ lại chi phí bước trước.\\
\file{demand_std} & $\sigma_D$ & 1 & Độ lệch chuẩn nhiễu nhu cầu.\\
\file{traffic_std} & $\sigma_T$ & 1 & Độ lệch chuẩn nhiễu giao thông.\\
\file{traffic_offset} & $m$ & 5 & Độ lệch cộng thêm vào chi phí.\\
\file{max_steps_factor} & --- & 5 & Giới hạn episode bằng $5M$ bước.\\
\bottomrule
\end{tabularx}
\end{table}
\file{EnvConfig} là dataclass bất biến. Hàm kiểm tra cấu hình yêu cầu
$1\leq N<n^2$, $C\geq6$, $0\leq\beta<1$, độ lệch chuẩn không âm và hệ
số giới hạn bước dương. Trong huấn luyện, số khách được lấy lại theo
từng episode nên không cố định bằng giá trị mặc định 20.

\section{Khởi tạo bài toán và quản lý ngẫu nhiên}
Hàm \file{reset()} đặt kho vào tọa độ
$(\lfloor n/2\rfloor,\lfloor n/2\rfloor)$ bằng cách hoán đổi tọa độ với
nút 0. Sau đó chương trình chọn $N$ vị trí khác kho không hoàn lại.
Khoảng cách nền tính từ tọa độ sau hoán đổi:
\begin{equation}
b_{ij}=|x_i-x_j|+|y_i-y_j|.
\end{equation}
Đối với lưới có cạnh chẵn, quy tắc này chọn một ô gần tâm hình học;
không giả định tồn tại một ô tâm duy nhất.

Một seed được tách thành bốn luồng sinh số ngẫu nhiên cho bố trí khách,
nhu cầu, giao thông và hành động heuristic. Nhu cầu thật được sinh một
lần theo từng khách và giữ kín. Cách này giúp các phương pháp nhận cùng
bài toán với cùng seed, không để thứ tự thăm làm đổi nhu cầu của khách.
Các thuật toán vẫn có thể trải qua số bước và chuỗi hành động khác nhau.

\section{Nhu cầu danh nghĩa và nhu cầu thực}
Với mỗi khách $i$, chương trình sinh:
\begin{align}
\widetilde d_i &\sim \mathcal U\{1,2,3,4,5\},\\
d_i &=\min\left(C-1,\max\left(1,
\operatorname{round}(Z_i)\right)\right),
\quad Z_i\sim\mathcal N(\widetilde d_i,\sigma_D^2).
\end{align}
Phép làm tròn dùng \file{np.rint}. Do làm tròn và chặn phân phối,
$\widetilde d_i$ là mức danh nghĩa, không nhất thiết bằng đúng $\E[d_i]$.
Mảng \file{_true} chỉ phục vụ simulator; mảng \file{revealed} khởi đầu
bằng NaN và được cập nhật khi xe đến khách.

Giá trị nhu cầu mà tác tử được biết tại bước $t$ là:
\begin{equation}
\widehat d_i^{(t)}=
\begin{cases}
\widetilde d_i,&\text{chưa tiết lộ nhu cầu},\\
d_i,&\text{đã tiết lộ nhu cầu}.
\end{cases}
\end{equation}
Hàm \file{known_demands()} thực hiện quy tắc này; nhu cầu của nút không
có khách được giữ bằng 0.

\section{Mô hình giao thông}
Hàm \file{_update_traffic()} tính:
\begin{equation}
c_{ij}^{\mathrm{new}}=
\max\left(b_{ij},\beta c_{ij}^{\mathrm{old}}
 +(1-\beta)b_{ij}+\varepsilon_{ij}+m\right),
\quad\varepsilon_{ij}\sim\mathcal N(0,\sigma_T^2).
\end{equation}
Sau đó đặt $c_{ii}=0$. Nhiễu theo các cạnh có hướng làm ma trận chi phí
có thể bất đối xứng dù khoảng cách nền đối xứng.

Môi trường cập nhật giao thông một lần khi reset và một lần trước mỗi
di chuyển. Vì vậy, chi phí thực trả cho hành động được lấy từ ma trận
mới, còn quan sát dùng để chọn hành động chứa ma trận trước cập nhật
đó. Đây là thứ tự cụ thể của mã, cần giữ nguyên khi so sánh thực nghiệm.

\section{Quan sát và mặt nạ hành động}
Hàm \file{observation()} trả ma trận $X_t\in\mathbb R^{M\times6}$.
Mỗi hàng là:
\begin{equation}
X_{t,i}=\left[
\widehat d_i^{(t)},\widetilde d_i,
\ind\{i\text{ đã phục vụ}\},\ind\{i=p_t\},
\ell_t,c_{p_t,i}^{(t)}\right].
\end{equation}
Với tải hữu hạn, $\ell_t=L_t$. Trong kịch bản sức chứa vô hạn, chỉ riêng
đặc trưng tải được mã hóa bằng 15 để tránh đưa Inf vào mạng; sức chứa
thật trong simulator vẫn là vô hạn.

Mask khách hàng chỉ cho phép chọn khách chưa phục vụ, không trùng vị
trí hiện tại và có nhu cầu đã biết không vượt tải. Kho được chọn khi xe
ở ngoài kho; khi đang ở kho và còn khách khả thi, hành động tự lặp tại
kho bị chặn. Nếu vừa giao thất bại, mask chỉ cho phép về kho.

\section{Chuyển trạng thái và xử lý thiếu hàng}\label{sec:recourse}
Hàm \file{step(action)} thực hiện lần lượt: kiểm tra episode và hành
động, cập nhật traffic, cộng chi phí, đổi vị trí, tiết lộ nhu cầu, xử lý
phục vụ và tính reward. Đoạn xử lý trung tâm là:
\begin{lstlisting}[caption={Xử lý nhu cầu thực khi xe đến khách}]
self.revealed[action] = self._true[action]
if self.revealed[action] <= self.load:
    self.load -= self.revealed[action]
    self.served[action] = True
    delivered = 1
else:
    self.failures += 1
    self.force_depot = True
\end{lstlisting}
Khi thiếu hàng, chuyến đi vẫn mất chi phí nhưng không trừ tải và không
đánh dấu khách đã phục vụ. Xe phải về kho để bổ sung hàng, sau đó khách
còn trong danh sách cần phục vụ. Như vậy, khách được phục vụ thành công
một lần nhưng có thể được thăm nhiều lần. Chương trình không bảo đảm
ràng buộc ``mỗi khách được thăm đúng một lần'' trong trường hợp này.

Reward mỗi bước được tính bằng:
\begin{equation}
r_t=-\Delta C_t+5\Delta N_t-0{,}5.
\end{equation}
Trong đó $\Delta N_t=1$ chỉ khi phục vụ thành công một khách mới.
Episode hoàn thành khi mọi khách đã phục vụ và xe về kho. Nếu đạt $5M$
bước mà chưa xong, episode được đánh dấu \file{truncated}.

\chapter{Kiến trúc mạng RoutingAgent}
\section{Đầu vào, đầu ra và kích thước}
Mạng nhận ma trận đặc trưng $X_t$, ma trận chi phí $C_t$ và bộ nhớ GRU
từ bước trước. Mạng trả logits $z_t\in\mathbb R^M$, giá trị vô hướng
$V_t$ và bộ nhớ mới $h_t\in\mathbb R^{1\times1\times64}$.

\begin{table}[htbp]
\centering\small
\caption{Kích thước các lớp của RoutingAgent}
\begin{tabularx}{\textwidth}{@{}lYl@{}}
\toprule
Thành phần & Cấu hình & Đầu ra \\
\midrule
NNConv 1 + ReLU & $6\rightarrow64$, mean aggregation & $M\times64$\\
NNConv 2 + ReLU & $64\rightarrow64$, mean aggregation & $M\times64$\\
Transformer & 2 lớp, 4 heads, FFN 128, dropout 0 & $M\times64$\\
Mean pooling & Trung bình trên toàn bộ nút & $1\times1\times64$\\
GRU & Input 64, hidden 64 & $1\times1\times64$\\
Fusion + tanh & Nối embedding nút và context, $128\rightarrow64$ & $M\times64$\\
Actor head & Linear $64\rightarrow1$ cho từng nút & $M$\\
Critic head & Linear $64\rightarrow64$, ReLU, Linear $64\rightarrow1$ & Vô hướng\\
\bottomrule
\end{tabularx}
\end{table}

\section{Thông điệp cạnh với NNConv}
Mỗi chi phí cạnh đi qua một MLP để sinh ma trận trọng số:
\begin{equation}
W_{ji}=\operatorname{reshape}\left(
\operatorname{MLP}_{\mathrm{edge}}(c_{ji}),d_{\mathrm{in}},d_{\mathrm{out}}
\right).
\end{equation}
Hai edge MLP có cấu trúc $1\rightarrow32\rightarrow6\times64$ và
$1\rightarrow32\rightarrow64\times64$. Nếu $x_j$ là vector hàng,
đầu ra lớp tại nút $i$ trước ReLU là:
\begin{equation}
y_i=x_iW_{\mathrm{root}}^{\mathsf T}+
\frac{1}{M-1}\sum_{j\ne i}x_jW_{ji}+b.
\end{equation}
Ma trận $W_{\mathrm{root}}$ có kích thước
$d_{\mathrm{out}}\times d_{\mathrm{in}}$, theo quy ước Linear trong
PyTorch. NNConv ở đây dùng trung bình thông điệp; cơ chế attention
đa đầu nằm trong Transformer, không có lớp GAT riêng.

\section{Tính toán NNConv theo nhóm cạnh}
Nếu lưới $15\times15$ thì $M=225$ và có $M(M-1)=50.400$ cạnh có hướng
không tự nối. Lớp thứ hai cần 4.096 số cho ma trận trọng số mỗi cạnh.
Materialize đồng thời các ma trận này cần khoảng 826 MB float32, chưa
tính bộ nhớ gradient.

Hàm \file{_conv()} chia cạnh thành các nhóm 256 cạnh. Với mỗi nhóm,
chương trình sinh trọng số, nhân thông điệp rồi dùng \file{index_add}
để cộng vào nút đích. Khi huấn luyện, activation checkpointing cho phép
tính lại một phần trung gian lúc backward, đổi thêm phép tính lấy giảm
bộ nhớ lưu trữ. Tập cạnh và phép trung bình được giữ nguyên. Bộ kiểm thử
so sánh cả output và gradient với NNConv chuẩn của PyTorch Geometric.

\section{Transformer, GRU và kết hợp bộ nhớ}
Sau Transformer, lấy trung bình embedding các nút:
\begin{align}
H_t&=\operatorname{Transformer}\left(H_t^{\mathrm{GNN}}\right),\\
g_t&=\frac1M\sum_{i=0}^{M-1}H_{t,i},\\
h_t&=\operatorname{GRU}(g_t,h_{t-1}).
\end{align}
Context chung $h_t$ được nối với từng embedding nút, sau đó biến đổi
phi tuyến:
\begin{align}
u_{t,i}&=\tanh\left(W_f[H_{t,i};h_t]+b_f\right),\\
z_{t,i}&=w_{\pi}^{\mathsf T}u_{t,i}+b_{\pi},\\
V_t&=\operatorname{MLP}_{V}(h_t).
\end{align}
Đây là lựa chọn fusion của bản dựng vì bài báo chưa cho công thức chi
tiết. Nếu chỉ cộng một context chung rồi đưa qua head tuyến tính, phần
đóng góp chung có thể bị softmax triệt tiêu. Fusion phi tuyến giúp bộ
nhớ có khả năng thay đổi tương quan xác suất giữa các nút.

\begin{lstlisting}[caption={Luồng forward sau hai lớp NNConv}]
x = self.transformer(x.unsqueeze(0)).squeeze(0)
pooled = x.mean(0).view(1, 1, -1)
_, memory = self.gru(pooled, memory)
context = memory[-1, 0]
fused = torch.tanh(self.fusion(
    torch.cat((x, context.expand_as(x)), -1)))
return (self.policy(fused).squeeze(-1),
        self.value(context).squeeze(), memory)
\end{lstlisting}

\section{Phân phối hành động}
Gọi $\mathcal A_t$ là tập hành động sau mask. Xác suất hành động là:
\begin{equation}
\pi_{\theta}(a\mid o_t,h_{t-1})=
\begin{cases}
\displaystyle\frac{\exp z_{t,a}}{\sum_{j\in\mathcal A_t}\exp z_{t,j}},
&a\in\mathcal A_t,\\
0,&a\notin\mathcal A_t.
\end{cases}
\end{equation}
Hàm \file{distribution()} gán logits hành động bị chặn bằng âm vô hạn
rồi tạo \file{Categorical}. Nếu mask không có hành động, hàm báo lỗi.
Mặc định suy luận lấy mẫu theo phân phối; tùy chọn
\file{--deterministic} chuyển sang argmax.

\chapter{Huấn luyện bằng PPO}
\section{Thu thập quỹ đạo}
Hàm \file{ppo_episode()} dùng chính sách hiện tại chạy đến khi hoàn
thành hoặc đạt giới hạn bước. Ở từng bước, chương trình lưu đặc trưng,
chi phí cạnh, mask, bộ nhớ đầu vào, hành động, log-probability cũ và giá
trị critic. Phần rollout chạy trong \file{torch.no_grad()} để không giữ
đồ thị đạo hàm cho toàn episode.

\section{Return và advantage}
Return có chiết khấu được tính ngược từ cuối quỹ đạo:
\begin{align}
G_t&=r_t+\gamma G_{t+1},\quad G_T=0,\quad\gamma=0{,}5,\\
A_t&=G_t-V_{\mathrm{old},t},\\
\widehat A_t&=\frac{A_t-\overline A}
{\max(\operatorname{std}_{\mathrm{pop}}(A),10^{-8})}.
\end{align}
Giới hạn bước được coi là horizon hữu hạn nên không bootstrap từ trạng
thái cuối. Chương trình không dùng Generalized Advantage Estimation.
Mạng value có hai lớp tuyến tính như mô tả của bài báo.

\section{Hàm mất mát}
Với xác suất hành động cũ đã được lưu, tỉ số PPO là:
\begin{equation}
\rho_t(\theta)=\exp\left(
\log\pi_{\theta}(a_t\mid o_t,h_{t-1})-
\log\pi_{\mathrm{old}}(a_t\mid o_t,h_{t-1})\right).
\end{equation}
Hàm mất mát mà chương trình tối thiểu hóa là:
\begin{align}
\mathcal L_{\mathrm{actor}}
&=-\frac1T\sum_t\min\left(
\rho_t\widehat A_t,
\operatorname{clip}(\rho_t,1-\epsilon,1+\epsilon)\widehat A_t\right),\\
\mathcal L
&=\mathcal L_{\mathrm{actor}}
+\frac{0{,}5}{T}\sum_t(V_t-G_t)^2
-\frac{0{,}01}{T}\sum_t\mathcal H(\pi_{\theta,t}),
\quad\epsilon=0{,}2.
\end{align}
Hệ số value 0,5 là lựa chọn bổ sung; entropy 0,01 tham khảo notebook.
Mỗi episode cập nhật bốn epochs bằng Adam với learning rate $10^{-4}$.
Loss từng bước được chia cho độ dài quỹ đạo, tích lũy gradient rồi mới
clip chuẩn gradient về tối đa 1 và thực hiện optimizer step.

\begin{lstlisting}[caption={Tính loss và cập nhật tham số trong một epoch PPO}]
ratio = (dist.log_prob(action) - old_logp).exp()
surrogate = torch.minimum(
    ratio * advantages[idx],
    ratio.clamp(1-clip, 1+clip) * advantages[idx])
loss = (-surrogate + .5*(value-returns[idx]).square()
        - .01*dist.entropy()) / len(records)
loss.backward()
# After accumulating all trajectory steps:
torch.nn.utils.clip_grad_norm_(
    agent.parameters(), 1., error_if_nonfinite=True)
optimizer.step()
\end{lstlisting}

\section{Xử lý bộ nhớ khi replay}
Hidden state được lưu từ rollout và detach giữa các bước. Khi cập nhật
PPO, mỗi bước replay dùng hidden đầu vào đã lưu; không lan truyền
gradient qua toàn bộ chuỗi thời gian. Vì vậy đây là huấn luyện recurrent
với gradient horizon một bước. GRU vẫn nhận gradient từ đầu ra hiện tại,
nhưng không tương đương backpropagation xuyên suốt episode.

Dropout của Transformer được đặt bằng 0 để tránh thay đổi likelihood do
mask dropout giữa rollout và replay. Thiết lập này khác dropout mặc định
trong notebook tham chiếu và phải được công khai khi báo cáo kết quả.

\section{Lịch huấn luyện và checkpoint}
CLI mặc định đặt 200.000 episode, lưới $8\times8$, mỗi episode lấy đều
số khách từ 10 đến 63. Bài báo có chỗ ghi tối đa 50 khách; tham số
\file{--max-customers 50} cho phép khảo sát cách diễn giải đó.

Mỗi episode được gắn seed từ seed tổng và chỉ số episode. Checkpoint
lưu trọng số mô hình, optimizer, số episode, cấu hình và trạng thái RNG
Torch/CUDA. File tạm được ghi xong trước khi thay checkpoint chính.
Khi resume, chương trình kiểm tra lưới, seed và khoảng số khách có khớp
checkpoint hay không. Log dạng JSONL lưu thêm tuyến, số khách phục vụ và
số lần thiếu hàng để kiểm tra quá trình học.

\chapter{Các phương pháp giải và thiết kế đánh giá}
\section{Ba biến thể chính sách}
Ba biến thể dùng chung một mạng đã huấn luyện; khác biệt nằm ở mask
trong giai đoạn suy luận:
\begin{description}
\item[GNN--DRL (\file{gnn})] Dùng toàn bộ tập hành động khả thi.
\item[GNN--DRL--G (\file{gnn_g})] Chỉ giữ tối đa ba khách khả thi có chi
phí hiện tại nhỏ nhất từ vị trí xe; giữ hành động về kho nếu được phép.
\item[GNN--DRL--CW (\file{gnn_cw})] Tính savings của các cặp khách khả thi:
\begin{equation}
s_{ij}=c_{0i}+c_{0j}-c_{ij}.
\end{equation}
Sắp cặp theo savings giảm dần, chọn tối đa ba ứng viên và kiểm tra tổng
nhu cầu của tập được chọn không vượt tải. Nếu không chọn được ứng viên,
dùng ba khách khả thi gần nhất. Sau đó lấy mẫu bằng xác suất chính sách
trên tập được giữ và kho nếu khả thi.
\end{description}
Công thức savings của hybrid bám notebook. Baseline CW dùng savings
có hướng $c_{i0}+c_{0j}-c_{ij}$; hai công thức không đồng nhất trong
môi trường bất đối xứng. Chương trình hiện không huấn luyện lại một
policy riêng cho từng mask lai.

\section{Bảy phương pháp đối chứng}
\begin{table}[htbp]
\centering\small
\caption{Quy tắc chọn hành động của các baseline trong mã}
\begin{tabularx}{\textwidth}{@{}p{3.1cm}Y@{}}
\toprule
Phương pháp & Quy tắc \\
\midrule
Greedy & Chọn khách khả thi có chi phí hiện tại nhỏ nhất.\\
Tabu & Chọn gần nhất ngoài danh sách tabu, tenure 20; fallback nếu cần.\\
Random & Chọn ngẫu nhiên đều trong khách khả thi bằng RNG riêng.\\
Largest-Demand & Chọn nhu cầu đã biết lớn nhất; khoảng cách phá hòa.\\
Max-Reachable & Tối đa hóa số khách còn vừa tải sau ứng viên; khoảng cách phá hòa.\\
Clarke--Wright & Tái dựng savings merges có hướng và ràng buộc tải ở mỗi quyết định; ưu tiên đầu tuyến dài rồi khoảng cách.\\
Sweep & Chọn theo góc cực so với kho trong tập khả thi; khoảng cách phá hòa.\\
\bottomrule
\end{tabularx}
\end{table}
Các baseline dùng cùng simulator, mask, quy tắc thiếu hàng và giới hạn
bước. Baseline Tabu là phiên bản chọn gần nhất có danh sách cấm theo mô
tả bài báo; không phải một thuật toán tìm kiếm lân cận toàn tuyến mạnh.

\section{Kịch bản và kích thước bài toán}
Các cặp $(N,n)$ là $(20,8)$, $(50,8)$, $(100,12)$ và $(200,15)$.
Mặc định mỗi cặp có 100 seed cho từng kịch bản.
\begin{table}[htbp]
\centering
\caption{Bảy cấu hình thí nghiệm}
\begin{tabularx}{\textwidth}{@{}lcccc@{}}
\toprule
Tên & $\sigma_D$ & $\sigma_T$ & $m$ & $C$\\
\midrule
\file{standard} & 1 & 1 & 5 & 15\\
\file{demand_only} & 1 & 0 & 0 & 15\\
\file{traffic_only} & 0 & 1 & 5 & 15\\
\file{infinite_capacity} & 1 & 1 & 5 & $\infty$\\
\file{high_traffic} & 1 & 10 & 5 & 15\\
\file{high_demand} & 3 & 1 & 5 & 15\\
\file{no_uncertainty} & 0 & 0 & 0 & 15\\
\bottomrule
\end{tabularx}
\end{table}
Ở kịch bản vô hạn, phép chặn trên $C-1$ của nhu cầu không còn giới hạn
14 như khi $C=15$. Các phương pháp vẫn nhận cùng instance cho cùng seed.

\section{Chỉ số và nguyên tắc tổng hợp}
Chi phí của một lượt là tổng chi phí mọi chặng đã thực hiện. Tỉ lệ hoàn
thành trên $K$ mẫu:
\begin{equation}
q=\frac1K\sum_{k=1}^K\ind\{\text{phục vụ đủ và trở về kho}\}.
\end{equation}
Mean cost và standard error chỉ tính trên tập lượt hoàn thành
$\mathcal C$:
\begin{equation}
\overline C=\frac1{|\mathcal C|}\sum_{k\in\mathcal C}C_k,
\qquad SE=\frac{s_C}{\sqrt{|\mathcal C|}}.
\end{equation}
Nếu có ít hơn hai lượt hoàn thành thì không báo SE; không có lượt hoàn
thành thì không báo mean cost. Khi các phương pháp có mức hoàn thành
khác nhau, cần đọc cost cùng completion rate; các trung bình trên tập
hoàn thành riêng không phải phép so sánh ghép cặp hoàn toàn.

Z-score và Friedman chỉ dùng những seed mà tất cả phương pháp đang
so sánh đều hoàn thành. Với mỗi seed hợp lệ, chuẩn hóa chi phí qua các
phương pháp:
\begin{equation}
z_{k,m}=\frac{C_{k,m}-\mu_k}{\max(\sigma_k,10^{-12})}.
\end{equation}
Sau đó lấy trung bình $z_{k,m}$ theo seed. Friedman được chạy riêng từng
kịch bản và kích thước, cần tối thiểu ba phương pháp và hai block không
hoàn toàn đồng hạng. Đây là kiểm định tổng thể; p-value nhỏ chưa chứng
minh mô hình thắng riêng từng đối thủ. Chưa cài kiểm định hậu nghiệm.

Runtime đo sau khi instance đã được tạo, gồm bước dựng đặc trưng, giải
và cập nhật môi trường. Khi dùng CUDA, chương trình đồng bộ trước và
sau phép đo. Khác phần cứng và quy tắc timing có thể tạo khác biệt lớn
với số liệu runtime của bài báo.

\section{Tệp đầu ra}
\begin{table}[htbp]
\centering\small
\caption{Dữ liệu do bộ đánh giá xuất ra}
\begin{tabularx}{\textwidth}{@{}p{3.5cm}Y@{}}
\toprule
Tệp & Nội dung \\
\midrule
\file{instances.csv} & Cost, runtime, served, complete, truncated, steps, failures từng lượt.\\
\file{summary.csv} & Mean cost, SE, completion rate, runtime, failures, paired mean z-score.\\
\file{friedman.json} & Thống kê, p-value, số block hợp lệ hoặc lý do chưa kiểm định.\\
\file{routes.jsonl} & Tuyến, chi phí từng chặng, tọa độ và tập khách.\\
\file{settings.json} & Tham số chạy; episode, SHA-256 và phiên bản implementation của checkpoint nếu dùng mạng.\\
\bottomrule
\end{tabularx}
\end{table}

\chapter{Kiểm thử, kết quả vận hành và giới hạn}
\section{Các nhóm kiểm thử}
Theo \file{VALIDATION.md}, lần kiểm tra triển khai ghi nhận
\textbf{69 kiểm thử đạt}. Các nhóm gồm:
\begin{enumerate}
\item Kiểm tra tọa độ kho, Manhattan distance, seed và không lộ nhu cầu thật.
\item Kiểm tra công thức traffic, mask, hành động không hợp lệ và giới hạn bước.
\item Kiểm tra giao thất bại, quay về kho, phục vụ lại và tổng chi phí từng chặng.
\item Chạy bảy baseline trên bảy kịch bản.
\item So sánh đầu ra và gradient NNConv chia nhóm với NNConv chuẩn.
\item Kiểm tra GRU ảnh hưởng xác suất hành động và nhận gradient.
\item Kiểm tra discounted returns, cập nhật actor/critic, lưu/nạp trọng số.
\item Kiểm tra rollouts chính sách, xử lý dữ liệu chưa hoàn thành và thống kê.
\item Kiểm tra forward ở lưới 8, 12, 15 với 50, 100, 200 khách.
\end{enumerate}

\section{Kết quả chạy thử được ghi nhận}
Môi trường kiểm tra là Windows, Python 3.12, PyTorch 2.6.0 CPU và
PyTorch Geometric 2.6.1. Hai episode thử đầu tiên trên lưới $8\times8$
với 10 khách đều phục vụ 10/10 và kết thúc ở kho. Sau khi resume,
episode thứ ba cũng hoàn thành 10/10 khách.

Bộ đánh giá checkpoint đã chạy đủ mười phương pháp trên hai mẫu chuẩn
20 khách, ghi nhận 20/20 lượt hoàn thành. Riêng kiểm tra baseline gồm
$7\times7\times3=147$ lượt với 20 khách, ghi nhận 147/147 hoàn thành.
Đây là các kết quả chạy thử đã được lưu trong biên bản triển khai,
không phải một bộ thực nghiệm đủ lớn để xếp hạng chất lượng thuật toán.

Kiểm thử forward ở 200 khách xác nhận mạng xử lý được kích thước đầu
vào đó; không chứng minh chính sách đã học có chất lượng tốt hoặc suy
luận đủ nhanh trên toàn tuyến 200 khách.

\section{Những khác biệt cần công khai}
\begin{longtable}{@{}p{3.3cm}p{11.2cm}@{}}
\caption{Các lựa chọn triển khai ảnh hưởng khả năng tái lập}\label{tab:limits}\\
\toprule
Nội dung & Cách triển khai và hệ quả\\
\midrule
\endfirsthead
\toprule Nội dung & Cách triển khai và hệ quả\\\midrule
\endhead
\bottomrule
\endfoot
Số khách huấn luyện & Mặc định 10--63 theo mục phương pháp; có tùy chọn 50 theo phần kết luận. Không xem hai cấu hình là cùng thí nghiệm.\\
Biểu diễn đồ thị & Dùng toàn bộ $n^2$ ô theo notebook thay vì chỉ $N+1$ nút; có thể ảnh hưởng pooling và attention.\\
Critic & Dùng mạng value theo PDF, trong khi notebook tham chiếu dùng mean return làm baseline.\\
GRU và fusion & Dùng hidden GRU trong actor/critic và kết hợp phi tuyến; là chi tiết bổ sung do bài chưa nêu công thức đầy đủ.\\
Thiếu tải & Có recourse quay kho và thăm lại; không giữ được ràng buộc thăm đúng một lần trong mọi tình huống.\\
Recurrent replay & Hidden được detach và replay từ rollout; gradient horizon một bước, không phải BPTT toàn chuỗi.\\
Hybrid CW & Mask savings được tính lại và vẫn lấy mẫu bằng policy; không chạy giống hệt nhánh gần như tất định của notebook gốc.\\
Baseline & Dùng chung mask, traffic update và horizon; cách thống nhất này khác một số evaluator trong notebook.\\
Seed & Nhu cầu ẩn được gắn theo khách để ghép cặp; không sinh theo thứ tự thăm như notebook.\\
Dropout & Đặt bằng 0 để rollout/replay nhất quán; khác mặc định Transformer trong notebook.\\
\end{longtable}

\section{Giới hạn hiện tại}
Chương trình chưa được huấn luyện đủ 200.000 episode, chưa đánh giá toàn
bộ 100 mẫu trên bốn kích thước và chưa khớp cost, SE, p-value của các
bảng trong bài. Chưa có checkpoint và test instances gốc để đối chiếu.
Chưa xác minh trên Tesla T4, dữ liệu đường thực hoặc đơn hàng thực.

Môi trường dùng traffic nhiễu Gaussian đơn giản; không mô phỏng đầy đủ
giờ cao điểm và tương quan tắc nghẽn không gian. Giá trị $\gamma=0{,}5$
được giữ theo bài nhưng làm trọng số reward xa giảm nhanh. Các vấn đề
này cần được nghiên cứu bằng thực nghiệm, không thể giải quyết chỉ bằng
việc chương trình chạy hết kiểm thử.

\chapter{Hướng dẫn sử dụng và hướng phát triển}
\section{Cài đặt và kiểm thử}
Các lệnh dưới đây chạy tại thư mục \file{research/cvrp_reproduction}.
Dấu gạch chéo trong đường dẫn được PowerShell trên Windows hỗ trợ.
\begin{lstlisting}[language={},caption={Tạo môi trường Python và cài thư viện trên CPU}]
python -m venv .venv
.venv/Scripts/python.exe -m pip install torch==2.6.0 --index-url https://download.pytorch.org/whl/cpu
.venv/Scripts/python.exe -m pip install -r requirements.txt
.venv/Scripts/python.exe -m pytest -q
\end{lstlisting}
Máy NVIDIA tương thích CUDA 12.4 có thể thay nguồn wheel Torch CPU bằng
\url{https://download.pytorch.org/whl/cu124}. Việc CUDA có chạy được trên
máy đích cần kiểm tra riêng.

\section{Huấn luyện thử và huấn luyện đầy đủ}
\begin{lstlisting}[language={},caption={Kiểm tra vận hành bằng hai episode}]
.venv/Scripts/python.exe -m cvrp.train --episodes 2 --min-customers 10 --max-customers 10 --output runs/smoke --save-every 1
\end{lstlisting}
\begin{lstlisting}[language={},caption={Huấn luyện theo cấu hình nghiên cứu trên GPU}]
.venv/Scripts/python.exe -m cvrp.train --episodes 200000 --device cuda --output runs/paper
\end{lstlisting}
Nếu không có GPU, bỏ \file{--device cuda} để tự chọn CPU. Thời gian
chạy dài phụ thuộc phần cứng; không có cam kết thời gian giống bài báo.

\begin{lstlisting}[language={},caption={Tiếp tục checkpoint với cùng cấu hình}]
.venv/Scripts/python.exe -m cvrp.train --resume runs/paper/checkpoint.pt --episodes 200000 --device cuda --output runs/paper
\end{lstlisting}
Tham số episodes là mốc tổng cần đạt, không phải số episode cộng thêm.
Chỉ nạp checkpoint do mình tạo hoặc đã tin cậy vì CLI hiện dùng định
dạng Torch checkpoint có lưu cả optimizer và metadata.

\section{Chạy đánh giá}
\begin{lstlisting}[language={},caption={Đánh giá thử checkpoint trên kịch bản chuẩn}]
.venv/Scripts/python.exe -m cvrp.evaluate --checkpoint runs/smoke/checkpoint.pt --instances 3 --sizes 20 --scenarios standard --methods gnn gnn_g gnn_cw tabu greedy random largest_demand max_reachable cw sweep --output runs/smoke_evaluation
\end{lstlisting}
\begin{lstlisting}[language={},caption={Đánh giá đủ bảy kịch bản và bốn kích thước}]
.venv/Scripts/python.exe -m cvrp.evaluate --checkpoint runs/paper/checkpoint.pt --device cuda --instances 100 --methods gnn gnn_g gnn_cw tabu greedy random largest_demand max_reachable cw sweep --output runs/full_evaluation
\end{lstlisting}
Lệnh đánh giá mặc định chỉ chạy baseline nếu không chỉ định methods.
Các phương pháp học bắt buộc có checkpoint; chương trình không âm thầm
dùng mạng khởi tạo ngẫu nhiên khi người dùng quên cung cấp trọng số.

\section{Hướng phát triển}
Các bước tiếp theo ưu tiên huấn luyện đầy đủ với nhiều training seed,
báo cáo completion rate cùng chi phí và đối chiếu baseline dưới cùng
ngân sách tính toán. Cần làm ablation cho GRU, Transformer, fusion và
mask CW để tách đóng góp của học và heuristic. Có thể khảo sát thêm
chiết khấu, reward, quy tắc recourse và recurrent replay dài hơn.

Để tiến gần ứng dụng, cần thay lưới tổng hợp bằng mạng đường có dữ liệu,
mô hình hóa traffic có tương quan và thử chuyển phân phối nhu cầu.
Bài toán nhiều xe, time windows và khách xuất hiện theo thời gian là
các phần mở rộng chưa được triển khai trong phiên bản hiện tại.

\chapter*{Kết luận}
\addcontentsline{toc}{chapter}{Kết luận}
Bản dựng cung cấp đầy đủ các thành phần phần mềm để khảo sát định tuyến
với bất định nhu cầu và giao thông: môi trường, NNConv--Transformer--GRU,
PPO, baseline, checkpoint và bộ đánh giá. Mã nguồn có các kiểm thử cho
logic môi trường, phép tính mạng và quy trình học.

Giá trị hiện được xác nhận là khả năng vận hành và kiểm tra có hệ thống.
Hiệu quả tối ưu sau huấn luyện dài và mức độ khớp bài báo vẫn cần được
xác minh bằng bộ thí nghiệm đầy đủ, với các lựa chọn triển khai được
công khai như trong báo cáo này.

\begin{thebibliography}{9}
\addcontentsline{toc}{chapter}{Tài liệu tham khảo}
\bibitem{kadyrov2025}
S. Kadyrov, A. Azamov, Y. Abdumajitov và C. Turan,
``Deep reinforcement learning for dynamic vehicle routing with demand
and traffic uncertainty'', \emph{Operations Research Perspectives},
tập 15, bài 100351, 2025.
\url{https://doi.org/10.1016/j.orp.2025.100351}.

\bibitem{notebook}
S. Kadyrov, \emph{CVRP Uncertainty}, notebook tham chiếu được dẫn ở
tài liệu [38] của bài báo.
\url{https://github.com/shiralikadyrov/CVRP_uncertainty}.
Git blob của bản tham chiếu:
\file{6f2e7293811039108151d2ef8f2979309155679c}.

\bibitem{implementation}
Mã triển khai cục bộ \file{research/cvrp_reproduction/cvrp/},
kèm \file{README.md}, \file{REPRODUCTION_NOTES.md} và
\file{VALIDATION.md}; biên bản kiểm tra ngày 15/11/2026.
\end{thebibliography}
\end{document}
